\section{Definitions}

\paragraph{Solipsis Objects}
The Solipsis world consists of a set of \emph{objects}. Each object is
associated with a unique identifier (UID) and may be classified into
groups: \emph{roaming} or \emph{static}. Roaming objects may roam
around the virtual world, and interact with it. Static objects on the
other hand are associated with a fixed position in the virtual
space. In this version of the specification, we consider only one type
of roaming and one type of static objects, respectively,
\emph{avatars} and \emph{sites}.

Objects may be described with several levels of details. In this
version of the specification, we concentrate on a very coarse-grained
subdivision between two such levels: a complete description,
comprising a full-fledged 3D model, and a simple predefined shape with
no textures that is encoded in an  \emph{object
  descriptor}

% Nonetheless, the protocol can
% easily support several levels of details with only minor
% modifications. The two levels we consider consist of a \emph{complete
%   description} of the object and of a an% \df{This will change. Ultimately we are going to have:
%   descriptor, and 3dData, which in turn consists of mesh, texture, and
%   animation.}.
The 3D description of the object contains all the information
necessary for the 3D visualization of an object, plus any sounds,
videos, or multimedia content associated with it. From the point of
view of this protocol, such a representation merely consists of a set
of binary objects, or files, that should be provided to the rendering
component. Each such object is associated with its own version number
and may thus be updated independently of the others.

The object descriptor, on the other hand, consist of the minimal
information required to display a rough representation of the object
and includes the object's location, its shape, its size, its
orientation, and its visibility radius. The shape of the object is
represented by an identifier that selects one from a set of predefined
shapes; while its visibility radius identifies the distance from which
the object should be visible in the absence of obstacles and without
any visual aids. 


The object descriptor of an object also keeps track of which hosts
have recorded a copy of the object's 3D description. Specifically,
each descriptor contains two fields associated with each of the files
that constitute the object's 3D model. The first field is the latest
known version number of the file, while the second is a list of hosts
that are known to have cached a copy of the version of the file
indicated in the first field.

The object descriptor is the means through which the \sol protocol
tracks the evolution of an object. To maintain consistency, each node
descriptor should be modified by only a single object at a time.  For
this reason, each descriptor also contains a sequence number that is
incremented each time the descriptor is modified.  The general rule
for modifying object descriptors in the \sol protocol is that each
descriptor may be modified only by the site node that is responsible
for the current location of the corresponding object. This is
straightforward in the case of static objects (sites), but is somewhat
more complex in the case of roaming objects (avatars), as these may
roam move from one location to another and thus from one site node to
another. 

Specifically, in the case of roaming object, the node responsible for
the object's location may not be the same as the node that is
currently owning the object. Yet, the owner of the object should still
be entitled to modify attributes of the object descriptor such as the
shape or the files associated with the 3D model. Details about how
this is achieved are given in Section~\ref{}. %  o make this
% possible, we also define for roaming objects, an additional
% \emph{static descriptor}\df{need a better name}, which includes all
% the fields in the object descriptors excluding location and
% orientation. The static descriptor of a roaming object can only be
% modified by the owner of the object, while the object descriptor is
% managed by the node responsible for the object location, as described
% in Section~\ref{}. 

Finally, in future versions of the protocol, the object descriptor may
contain information required to implement progressive levels of
details in the 3D representations of objects.
A summary of the fields included in an object descriptor is shown in
Table~\ref{tab:desc}.

\begin{table}
\begin{tabular}{|l|p{4cm}|}
UID & universal identifier of the object\\
owner & identifier of the node owning the object\\
Type & site or avatar\\
location & location in the 3D world\\
orientation & orientation in the 3D world\\
shape& shape from a predefined set\\
size & size of the object\\
visibility & distance from which the object is visible in the absence
of obstacles\\
\hline\\
seqNum & sequence number\\
\hline\\
$f_1$& first file for 3d model\\
$v_1$& version number for first file\\
$c_1$& list of hosts that have cached $v_1$ of $f_1$\\
...& ...\\
$f_n$& first file for 3d model\\
$v_n$& version number for first file\\
$c_n$& list of hosts that have cached $v_1$ of $f_1$\\
\hline\\
... & additional fields for progressive levels of details\\
\hline
\end{tabular}

\caption{Format of an object descriptor}
\label{tab:desc}
\end{table}

\paragraph{Solipsis Peers and Nodes}
A \emph{Solipsis peer}, or simply \emph{peer}, is any host that
participates in the \sol world by managing sites, avatars, or
both. Each peer is associated with a single instance of the \sol
platform, and throughout its lifetime, it may create or destroy its
objects according to the requests of its users. Moreover, it owns and
manages the objects it creates until they are
destroyed.\footnote{Future version of the specification should include
  the possibility to transfer the ownership of objects to other
  peers.} Finally, similar to an object, each peer is associated with
a unique identifier (UID).

Because each peer may be responsible for several objects at the same
time, we will use the term \emph{node} to identify the subset of
resources, and operations of a peer $p$, dedicated to the management
of a given object $o$. Specifically, we define a \sol \emph{node} as a
pair $(p, o)$ consisting of an object and of the corresponding
managing peer. Moreover, we will distinguish between \emph{site nodes}
and \emph{avatar nodes} according to what type of object the node
refers to.



\paragraph{3D Overlay}
Site nodes are organized in a three-dimensional overlay network whose
topology mimics the arrangement of the corresponding sites in the
virtual world.
% Nodes are organized in the overlay according to the RayNet
% protocol~\cite{} as described in Appendix~\ref{sec:raynet}.
Specifically, each node is associated with a location that is the same
as that of the corresponding site, and is responsible for a section of
the world, its \emph{cell}, centered around this location and
comprising a region of space around it. In simple terms, two nodes are
neighbors in the overlay if their corresponding sites are neighboring
in the \sol world. % This, mapping between overlay
% and 3D-world positions allows the  RayNet to provide efficient means to
% address and communicate with the nodes responsible for the sites
% encountered by avatars roaming in the \sol world.



%%% Local Variables: 
%%% mode: latex
%%% TeX-master: "specification"
%%% End: 
